\paragraph{Example 1 (qid 460) --- \rl{} 0.996.} Same \dieu{} as the reference: yes; same law: n/a; numbers grounded in sources: 100\%.
\begin{quote}\small\textbf{Question.} Theo Luật Các tổ chức tín dụng 2024, những hành vi nào bị cấm khi thực hiện hoạt động tín dụng?\\[2pt]
\textbf{Ground truth.} Căn cứ theo Điều 15 Luật Các tổ chức tín dụng 2024 quy định các hành vi bị nghiêm cấm trong hoạt động tín dụng như sau: - Tổ chức tín dụng, chi nhánh ngân hàng nước ngoài thực hiện hoạt động ngân hàng, hoạt động kinh doanh khác ngoài hoạt động ghi trong Giấy phép được Ngân hàng Nhà nước cấp cho tổ chức tín dụng, chi nhánh ngân hàng nước ngoài. - Tổ chức, cá nhân không phải là tổ chức tín dụng, chi nhánh ngân hàng nước ngoài thực hiện hoạt động ngân hàng, trừ giao dịch ký quỹ, giao dịch mua bán lại chứng khoán của …\\[2pt]
\textbf{System answer.} Căn cứ theo Điều 15 Luật Các tổ chức tín dụng 2024 quy định về hành vi bị nghiêm cấm trong hoạt động tín dụng như sau: - Tổ chức tín dụng, chi nhánh ngân hàng nước ngoài thực hiện hoạt động ngân hàng, hoạt động kinh doanh khác ngoài hoạt động ghi trong Giấy phép được Ngân hàng Nhà nước cấp cho tổ chức tín dụng, chi nhánh ngân hàng nước ngoài. - Tổ chức, cá nhân không phải là tổ chức tín dụng, chi nhánh ngân hàng nước ngoài thực hiện hoạt động ngân hàng, trừ giao dịch ký quỹ, giao dịch mua bán lại chứng khoán của …\end{quote}

\paragraph{Example 2 (qid 459) --- \rl{} 0.945.} Same \dieu{} as the reference: yes; same law: yes; numbers grounded in sources: 100\%.
\begin{quote}\small\textbf{Question.} Những tiêu chí và điều kiện cụ thể nào mà tổ chức tín dụng hoặc công ty Fintech phải thỏa để được cấp Giấy chứng nhận tham gia cơ chế thử nghiệm có kiểm soát trong lĩnh vực ngân hàng theo Nghị định 94/2025/NĐ-CP?\\[2pt]
\textbf{Ground truth.} Căn cứ theo Điều 8 Nghị định 94/2025/NĐ-CP quy định điều kiện và tiêu chí tham gia Cơ chế thử nghiệm có kiểm soát trong lĩnh vực ngân hàng như sau: - Tổ chức tín dụng không thuộc nhóm tổ chức tín dụng đang được kiểm soát đặc biệt theo Luật Các tổ chức tín dụng 2024, chi nhánh ngân hàng nước ngoài được xem xét cấp Giấy chứng nhận tham gia Cơ chế thử nghiệm khi giải pháp Fintech đáp ứng các tiêu chí như sau: + Là giải pháp có nội dung kỹ thuật và nghiệp vụ mà quy định pháp lý hiện hành chưa hướng dẫn cụ thể, rõ ràng …\\[2pt]
\textbf{System answer.} Căn cứ theo Điều 8 Nghị định 94/2025/NĐ-CP quy định về điều kiện và tiêu chí tham gia Cơ chế thử nghiệm có kiểm soát trong lĩnh vực ngân hàng như sau: - Tổ chức tín dụng không thuộc nhóm tổ chức tín dụng đang được kiểm soát đặc biệt được xem xét cấp Giấy chứng nhận tham gia Cơ chế thử nghiệm khi giải pháp Fintech đáp ứng các tiêu chí như sau: + Là giải pháp có nội dung kỹ thuật và nghiệp vụ mà quy định pháp luật hiện hành chưa hướng dẫn cụ thể, rõ ràng cho việc triển khai, áp dụng; + Là giải pháp có tính đổi mới …\end{quote}

\paragraph{Example 3 (qid 1865) --- \rl{} 0.899.} Same \dieu{} as the reference: yes; same law: n/a; numbers grounded in sources: 100\%.
\begin{quote}\small\textbf{Question.} Khi nộp hồ sơ xin cấp giấy phép thành lập và hoạt động cho một doanh nghiệp kinh doanh bảo hiểm, những tài liệu cần chuẩn bị theo Luật Kinh doanh bảo hiểm 2022 gồm những gì?\\[2pt]
\textbf{Ground truth.} Căn cứ theo Điều 69 Luật Kinh doanh bảo hiểm 2022, hồ sơ đề nghị cấp giấy phép thành lập và hoạt động kinh doanh bảo hiểm gồm các giấy tờ dưới đây:- Văn bản đề nghị cấp giấy phép thành lập và hoạt động.- Dự thảo điều lệ đối với doanh nghiệp bảo hiểm, doanh nghiệp tái bảo hiểm; dự thảo quy chế tổ chức và hoạt động đối với chi nhánh nước ngoài tại Việt Nam.- Phương án hoạt động 05 năm đầu, trong đó nêu rõ các nghiệp vụ bảo hiểm dự kiến triển khai, mô hình quản trị rủi ro, phương thức trích lập dự phòng nghiệp vụ, …\\[2pt]
\textbf{System answer.} Căn cứ theo Điều 69 Luật Kinh doanh bảo hiểm 2022 quy định về hồ sơ đề nghị cấp giấy phép thành lập và hoạt động như sau: - Văn bản đề nghị cấp giấy phép thành lập và hoạt động; - Dự thảo điều lệ đối với doanh nghiệp bảo hiểm, doanh nghiệp tái bảo hiểm; dự thảo quy chế tổ chức và hoạt động đối với chi nhánh nước ngoài tại Việt Nam; - Phương án hoạt động 5 năm đầu, trong đó nêu rõ các nghiệp vụ bảo hiểm dự kiến triển khai, mô hình quản trị rủi ro, phương thức trích lập dự phòng nghiệp vụ, chương trình tái bảo hiểm, …\end{quote}

\paragraph{Example 4 (qid 1136) --- \rl{} 0.820.} Same \dieu{} as the reference: yes; same law: yes; numbers grounded in sources: 100\%.
\begin{quote}\small\textbf{Question.} Từ 20/11/2024, khi tính mức lãi suất tối đa cho tiền gửi bằng VNĐ tại các tổ chức tín dụng thì có những khoản nào phải được tính vào mức tối đa đó?\\[2pt]
\textbf{Ground truth.} Căn cứ tại khoản 3 Điều 3 Thông tư 48/2024/TT-NHNN quy định về lãi suất như sau: Điều 3. Lãi suất 1. Tổ chức tín dụng áp dụng lãi suất tiền gửi bằng đồng Việt Nam của tổ chức và cá nhân không vượt quá mức lãi suất tối đa đối với tiền gửi không kỳ hạn, tiền gửi có kỳ hạn dưới 1 tháng, tiền gửi có kỳ hạn từ 1 tháng đến dưới 6 tháng do Thống đốc Ngân hàng Nhà nước Việt Nam quyết định trong từng thời kỳ và đối với từng loại hình tổ chức tín dụng. 2. Tổ chức tín dụng áp dụng lãi suất tiền gửi bằng đồng Việt Nam đối với …\\[2pt]
\textbf{System answer.} Căn cứ theo Điều 3 Thông tư 48/2024/TT-NHNN quy định về lãi suất như sau: Điều 3. Lãi suất 1. Tổ chức tín dụng áp dụng lãi suất tiền gửi bằng đồng Việt Nam của tổ chức và cá nhân không vượt quá mức lãi suất tối đa đối với tiền gửi không kỳ hạn, tiền gửi có kỳ hạn dưới 1 tháng, tiền gửi có kỳ hạn từ 1 tháng đến dưới 6 tháng do Thống đốc Ngân hàng Nhà nước Việt Nam quyết định trong từng thời kỳ và đối với từng loại hình tổ chức tín dụng. 2. Tổ chức tín dụng áp dụng lãi suất tiền gửi bằng đồng Việt Nam đối với tiền …\end{quote}

\paragraph{Example 5 (qid 2213) --- \rl{} 0.751.} Same \dieu{} as the reference: yes; same law: yes; numbers grounded in sources: 100\%.
\begin{quote}\small\textbf{Question.} Vần seri trên tờ tiền được hiểu như thế nào theo quy định hiện hành?\\[2pt]
\textbf{Ground truth.} Căn cứ khoản 1 Điều 3 Thông tư 01/2024/TT-NHNN quy định như sau:Điều 3. Giải thích từ ngữTrong Thông tư này, các từ ngữ dưới đây được hiểu như sau:1. Vần seri là phần chữ được ghép bởi 2 trong số 26 chữ cái (gồm: A, B, C, D, E, F, G, H, I, J, K, L, M, N, O, P, Q, R, S, T, U, V, W, X, Y, Z). Vần seri bao gồm vần chính và vần phụ.2. Vần chính là vần seri sử dụng để in trên các tờ tiền theo số lượng, cơ cấu tiền in hàng năm.3. Vần phụ là vần seri sử dụng để in trên các tờ tiền thay thế những tờ tiền in hỏng.4. Seri …\\[2pt]
\textbf{System answer.} Căn cứ tại khoản 2 Điều 3 Thông tư 01/2024/TT-NHNN quy định về vần seri như sau: Điều 3. Giải thích từ ngữ Trong Thông tư này, các từ ngữ dưới đây được hiểu như sau: 1. Vần seri là phần chữ được ghép bởi 2 trong số 26 chữ cái (gồm: A, B, C, D, E, F, G, H, I, J, K, L, M, N, O, P, Q, R, S, T, U, V, W, X, Y, Z). Vần seri bao gồm vần chính và vần phụ. 2. Vần chính là vần seri sử dụng để in trên các tờ tiền theo số lượng, cơ cấu tiền in hàng năm. 3. Vần phụ là vần seri sử dụng để in trên các tờ tiền thay thế những tờ …\end{quote}

\paragraph{Example 6 (qid 1941) --- \rl{} 0.677.} Same \dieu{} as the reference: yes; same law: yes; numbers grounded in sources: 100\%.
\begin{quote}\small\textbf{Question.} Theo Thông tư 20/2024, kể từ 01/7/2024, những loại tiền tệ nào được sử dụng để thực hiện bao thanh toán, trả nợ và có quy định gì về mua bán ngoại tệ khi khách hàng không có đủ ngoại tệ?\\[2pt]
\textbf{Ground truth.} Căn cứ theo Điều 9 Thông tư 20/2024/TT-NHNN quy định đồng tiền bao thanh toán, trả nợ của tổ chức tín dụng như sau: - Đối với bao thanh toán bên bán hàng, đồng tiền bao thanh toán là đồng tiền của khoản phải thu hoặc đồng Việt Nam. - Đối với bao thanh toán bên mua hàng: + Đồng tiền bao thanh toán là đồng Việt Nam, trừ trường hợp quy định tại điểm b khoản 2 Điều 9 Thông tư 20/2024/TT-NHNN; + Đơn vị bao thanh toán được xem xét, quyết định bao thanh toán bằng ngoại tệ đối với khoản phải thu bằng ngoại tệ phù hợp với …\\[2pt]
\textbf{System answer.} Căn cứ theo Điều 9 Thông tư 20/2024/TT-NHNN quy định về đồng tiền bao thanh toán, trả nợ như sau: Điều 9. Đồng tiền bao thanh toán, trả nợ 1. Đối với bao thanh toán bên bán hàng, đồng tiền bao thanh toán là đồng tiền của khoản phải thu hoặc đồng Việt Nam. 2. Đối với bao thanh toán bên mua hàng: a) Đồng tiền bao thanh toán là đồng Việt Nam, trừ trường hợp quy định tại điểm b khoản này; b) Đơn vị bao thanh toán được xem xét, quyết định bao thanh toán bằng ngoại tệ đối với khoản phải thu bằng ngoại tệ phù hợp với quy …\end{quote}

\paragraph{Example 7 (qid 2178) --- \rl{} 0.627.} Same \dieu{} as the reference: yes; same law: yes; numbers grounded in sources: 100\%.
\begin{quote}\small\textbf{Question.} Những quy định cụ thể về quản lý tài sản phần mềm trong hệ thống thông tin ngân hàng được Thông tư 09/2020/TT-NHNN quy định như thế nào?\\[2pt]
\textbf{Ground truth.} Căn cứ quy định Điều 10 Thông tư 09/2020/TT-NHNN quy định về quản lý tài sản phần mềm như sau: Quản lý tài sản phần mềm 1. Với mỗi hệ thống thông tin do tổ chức quản lý trực tiếp, tổ chức phải lập danh sách tài sản phần mềm với các thông tin cơ bản gồm: tên tài sản, giá trị, mục đích sử dụng, phạm vi sử dụng, chủ thể quản lý, thông tin về bản quyền, phiên bản, hệ thống thông tin thành phần (nếu có). 2. Tài sản phần mềm phải được gắn trách nhiệm cho cá nhân hoặc bộ phận quản lý. 3. Tài sản phần mềm phải được tổ …\\[2pt]
\textbf{System answer.} Căn cứ theo Điều 10 Thông tư 09/2020/TT-NHNN quy định về quản lý tài sản phần mềm như sau: - Với mỗi hệ thống thông tin do tổ chức quản lý trực tiếp phải lập danh sách tài sản phần mềm với các thông tin cơ bản gồm: tên tài sản, giá trị, mục đích sử dụng, phạm vi sử dụng, chủ thể quản lý, thông tin về bản quyền, phiên bản, hệ thống thông tin thành phần (nếu có). - Tài sản phần mềm phải được gắn trách nhiệm cho cá nhân hoặc bộ phận quản lý. - Tài sản phần mềm phải được tổ chức định kỳ rà soát và cập nhật các bản vá …\end{quote}

\paragraph{Example 8 (qid 2366) --- \rl{} 0.582.} Same \dieu{} as the reference: yes; same law: yes; numbers grounded in sources: 100\%.
\begin{quote}\small\textbf{Question.} Sổ/tấm thẻ tiết kiệm tại ngân hàng cần ghi những thông tin bắt buộc nào theo quy định?\\[2pt]
\textbf{Ground truth.} Căn cứ tại khoản 2 Điều 7 Thông tư 48/2018/TT-NHNN quy định nội dung thẻ tiết kiệm gồm: - Tên tổ chức tín dụng, con dấu; Họ tên, chữ ký của giao dịch viên và của người đại diện hợp pháp của tổ chức tín dụng; - Họ tên, số và ngày cấp Giấy tờ xác minh thông tin của người gửi tiền hoặc của tất cả người gửi tiền (đối với tiền gửi tiết kiệm chung) và thông tin của người đại diện theo pháp luật của người gửi tiền trong trường hợp gửi tiền gửi tiết kiệm thông qua người đại diện theo pháp luật; - Số Thẻ tiết kiệm; số …\\[2pt]
\textbf{System answer.} Căn cứ theo khoản 2 Điều 7 Thông tư 48/2018/TT-NHNN quy định về thẻ tiết kiệm như sau: Điều 7. Thẻ tiết kiệm 1. Thẻ tiết kiệm hoặc Sổ tiết kiệm (sau đây gọi chung là Thẻ tiết kiệm) là chứng chỉ xác nhận quyền sở hữu tiền gửi tiết kiệm của người gửi tiền tại tổ chức tín dụng, được áp dụng đối với trường hợp nhận tiền gửi tiết kiệm tại địa điểm giao dịch hợp pháp thuộc mạng lưới hoạt động của tổ chức tín dụng. 2. Nội dung Thẻ tiết kiệm a) Thẻ tiết kiệm phải có tối thiểu các nội dung sau: (i) Tên tổ chức tín dụng, …\end{quote}

\paragraph{Example 9 (qid 2339) --- \rl{} 0.501.} Same \dieu{} as the reference: no; same law: yes; numbers grounded in sources: 78\%.
\begin{quote}\small\textbf{Question.} Từ ngày 14/05/2024, quy trình quản lý và giao nhận các tờ tiền mới in được thực hiện như thế nào, bao gồm việc lập bảng kê seri khi chuyển giao giữa các đơn vị? Và cấu trúc mã seri trên các tờ tiền mới in được xây dựng ra sao (đặc biệt phân biệt tiền in trước và sau năm 2003)?\\[2pt]
\textbf{Ground truth.} Quy định quản lý seri tiền mới in giao nhận tiền mới in từ 14/05/2024 như sau: Căn cứ theo Điều 6 Thông tư 01/2024/TT-NHNN (Có hiệu lực từ 14/05/2024) quy định về quản lý seri tiền mới in giao nhận tiền mới in từ 14/05/2024 như sau: - Cơ sở in, đúc tiền giao tiền mới in cho Ngân hàng Nhà nước (Cục Phát hành và Kho quỹ); giao, nhận tiền mới in giữa các Kho tiền Trung ương; giao, nhận tiền mới in giữa Kho tiền Trung ương với Sở Giao dịch và Ngân hàng Nhà nước chi nhánh; giao, nhận tiền mới in giữa Ngân hàng Nhà nước …\\[2pt]
\textbf{System answer.} Căn cứ theo Điều 6 Thông tư 01/2024/TT-NHNN quy định về quản lý seri tiền mới in trong quá trình giao, nhận tiền như sau: - Trong quá trình giao, nhận tiền mới in + Cơ sở in, đúc tiền giao tiền mới in cho Ngân hàng Nhà nước (Cục Phát hành và Kho quỹ); giao, nhận tiền mới in giữa các Kho tiền Trung ương; giao, nhận tiền mới in giữa Kho tiền Trung ương với Sở Giao dịch và Ngân hàng Nhà nước chi nhánh; giao, nhận tiền mới in giữa Ngân hàng Nhà nước chi nhánh với nhau. + Bên giao tiền mới in phải lập bảng kê seri của …\end{quote}

\paragraph{Example 10 (qid 2412) --- \rl{} 0.408.} Same \dieu{} as the reference: no; same law: yes; numbers grounded in sources: 82\%.
\begin{quote}\small\textbf{Question.} Đồng 500 đồng có còn là phương tiện thanh toán hợp pháp tại Việt Nam hiện nay không? Trong giao dịch, người bán hoặc người nhận có được từ chối chấp nhận tờ tiền 500 đồng nếu tờ tiền đó đủ tiêu chuẩn lưu thông không?\\[2pt]
\textbf{Ground truth.} Tờ 500 đồng hiện nay có còn sử dụng hay không?Theo Thông báo 6099/NHNN-PHKQ năm 2007 thông báo về việc đình chỉ lưu hành tiền cotton loại 50.000 đồng và 100.000 đồng cụ thể như sau:1. Từ ngày 01/9/2007 đồng tiền cotton loại 50.000 đồng và 100.000 đồng hết giá trị lưu hành trên lãnh thổ Việt Nam.2. Kể từ thời điểm đình chỉ lưu hành (01/9/2007), các tổ chức và cá nhân có các loại tiền này được đổi ngang giá trị ra các loại tiền do Ngân hàng Nhà nước Việt Nam phát hành tại Ngân hàng Nhà nước chi nhánh tỉnh, thành …\\[2pt]
\textbf{System answer.} Căn cứ theo Điều 1 Thông tư 13/2010/TT-NHNN được sửa đổi bởi khoản 1 Điều 1 Thông tư 14/2012/TT-NHNN quy định về tiền Việt Nam như sau: Điều 1. Tiền Việt Nam 1. Tiền Việt Nam là phương tiện thanh toán hợp pháp của nước Cộng hoà xã hội chủ nghĩa Việt Nam. 2. Tiền Việt Nam bao gồm tiền giấy và tiền kim loại do Ngân hàng Nhà nước Việt Nam phát hành. Căn cứ theo Điều 4 Thông tư 25/2013/TT-NHNN quy định về tiêu chuẩn tiền không đủ tiêu chuẩn lưu thông như sau: Điều 4. Tiêu chuẩn tiền không đủ tiêu chuẩn lưu thông 1. …\end{quote}

\paragraph{Example 11 (qid 2085) --- \rl{} 0.288.} Same \dieu{} as the reference: n/a; same law: n/a; numbers grounded in sources: 33\%.
\begin{quote}\small\textbf{Question.} Tờ tiền mệnh giá 200.000 đồng in trên nền polymer được phát hành lần đầu vào thời điểm nào và những đặc điểm nhận dạng chính của tờ tiền đó là gì?\\[2pt]
\textbf{Ground truth.} Tờ 200.000 đồng Polymer phát hành vào năm nào? Tại Thông báo 7001/TB-NHNN năm 2006 quy định về việc phát hành tờ 200.000 đồng Polymer và tờ 10.000 đồng Polymer có quy định như sau: Thực hiện Quyết định số 1048/QĐ-TTg ngày 10/8/2006 của Thủ tướng Chính phủ, từ ngày 30/8/2006, Ngân hàng Nhà nước Việt Nam phát hành vào lưu thông đồng tiền mệnh giá 200.000 đồng và 10.000 đồng mới, in trên giấy nền Polymer. Như vậy, tờ 200.000 đồng Polymer phát hành vào ngày 30/8/2006. Tờ 200.000 đồng Polymer có những đặc điểm chính …\\[2pt]
\textbf{System answer.} Căn cứ theo Thông tư 11/2021/TT-NHNN quy định về tiền mới in của Ngân hàng Nhà nước Việt Nam như sau: Chương 2. Quy định cụ thể 2.1. Nguyên tắc quản lý ... 2.2. Quy định về phát hành, thu hồi và tiêu hủy tiền mới in ... 2.3. Quy định về quản lý seri tiền mới in Tại đây Theo đó, từ ngày 1/7/2021, Ngân hàng Nhà nước Việt Nam bắt đầu phát hành loại tiền mới in mệnh giá 200.000 đồng trên nền polymer. Đồng tiền mới in mệnh giá 200.000 đồng có các đặc điểm nhận dạng sau: - Vần seri: 200.000. - Số thứ tự bao: 001. - Năm …\end{quote}

\paragraph{Example 12 (qid 572) --- \rl{} 0.063.} Same \dieu{} as the reference: n/a; same law: n/a; numbers grounded in sources: 25\%.
\begin{quote}\small\textbf{Question.} Vietcombank thu mức phí thường niên hàng năm cho thẻ ghi nợ cá nhân ở khoảng bao nhiêu tiền?\\[2pt]
\textbf{Ground truth.} Phí thường niên là các khoản phí đóng định kỳ hàng năm để duy trì quyền lợi và sử dụng dịch vụ của khách hàng. Việc đóng phí thường niên sẽ giúp khách hàng tiếp tục được hưởng các quyền lợi từ dịch vụ mà mình sử dụng. Căn cứ biểu phí dịch vụ thẻ ghi nợ áp dụng cho khách hàng cá nhân tại Vietcombank quy định về mức phí thường niên như sau:Tải về: Biểu phí dịch vụ thẻ Ngân hàng Vietcombank Như vậy, phí thường niên Ngân hàng Vietcombank đối với thẻ ghi nợ áp dụng cho khách hàng cá nhân 1 năm giao động từ 54.545 đồng …\\[2pt]
\textbf{System answer.} Căn cứ theo Phụ lục 1 ban hành kèm theo Thông tư 35/2012/TT-NHNN quy định về phí dịch vụ thẻ ghi nợ nội địa như sau: Phí dịch vụ thẻ ghi nợ nội địa [1] Phí phát hành thẻ - Phí phát hành thẻ cho cá nhân: 100.000 đồng (một trăm triệu đồng) Việt Nam/thẻ. - Phí phát hành thẻ cho tổ chức: 1.000.000 đồng (một triệu đồng) Việt Nam/thẻ. [2] Phí giao dịch - Phí giao dịch rút tiền mặt: 100.000 đồng (một trăm nghìn đồng) Việt Nam/thẻ. - Phí giao dịch chuyển khoản: 10.000 đồng (mười nghìn đồng) Việt Nam/thẻ. - Phí giao dịch …\end{quote}
